\fbox{\parbox{0.92\textwidth}{\small
\textbf{Section XII: Mantra} (Chinese critical edition excerpt)\\[0.5em]
說般若波羅蜜多咒，卽說咒曰：揭帝揭帝，波羅揭帝，波羅僧揭帝，菩提僧莎訶\hspace{0pt}\textsuperscript{\textbf{a}}。\\[0.4em]
\rule{\linewidth}{0.3pt}\\[0.3em]
\textsuperscript{\textbf{a}}~說般若波羅蜜] 故說般若波羅蜜多呪，即說呪曰：揭諦揭諦，波羅揭諦，波羅僧揭諦，菩提薩婆訶。 \wit{Fangshan}; 故說般若波羅蜜呪，即說呪曰：竭帝竭帝，波羅竭帝，波羅僧竭帝，菩提僧莎呵。 \wit{T250}; 故說般若波羅蜜多咒。即說咒曰。唵 誐帝誐帝 波囉誐帝 波囉僧誐帝 菩提莎訶 \wit{T257}. The Fangshan reading 薩婆訶 (sàpóhē) is a more accurate transliteration of Sanskrit svāhā than T251's 僧莎訶 (sēng suōhē), which appears to be corrupt. 竭帝 (jiédì) vs T251 揭帝 (jiēdì) - different character for 'gate'
}}
